\sectionlink*{索引}

斜体のページ番号は演習を指す。

\begingroup
\small
\setlength{\parindent}{0pt}
\setlength{\parskip}{1pt}
\newcommand{\idx}[4]{\par\begingroup\leftskip=#1em\relax\hangindent=1em\relax\noindent#2\if\relax\detokenize{#3}\relax\else\space(#3)\fi\if\relax\detokenize{#4}\relax\else\quad#4\fi\par\endgroup}
% Original index: p.389
\idx{0}{Abel 微分}{Abelian differential}{293}
\idx{0}{作用（Dブレーン）}{action (D-brane)}{270–3}
\idx{1}{Born--Infeld 形式}{Born--Infeld form}{271, 273}
\idx{1}{線形化したもの}{linearized}{275–6}
\idx{0}{作用（一般）}{action (general)}{}
\idx{1}{Euclid 作用}{Euclidean}{335}
\idx{1}{フェルミオン振動子}{fermionic oscillator}{343}
\idx{1}{調和振動子}{harmonic oscillator}{334, 339–41}
\idx{1}{経路積分におけるもの}{in path integral}{77, 331, 343}
\idx{0}{作用（点粒子）}{action (point particle)}{9–11}
\idx{1}{Euclid 作用}{Euclidean}{82–3}
\idx{1}{ゲージ固定したもの}{gauge-fixed}{129}
\idx{1}{曲がった時空におけるもの}{in curved spacetime}{108}
\idx{1}{Wilson 線背景}{Wilson line background}{264}
\idx{0}{作用（時空）}{action (spacetime)}{}
\idx{1}{$c=1$ のモジュライ}{$c=1$ moduli}{246}
\idx{1}{閉弦タキオン}{closed string tachyon}{194}
\idx{1}{次元縮約}{dimensional reduction}{232–5, 249, 265}
\idx{1}{高階微分項}{higher-derivative terms}{188–9, 195}
\idx{1}{Hilbert 作用（重力）}{Hilbert (gravitational)}{83, 114, 194}
\idx{1}{継承原理}{inheritance principle}{258}
\idx{1}{局所対称性}{local symmetries}{196}
\idx{1}{質量ゼロの閉弦}{massless closed string}{114–15}
\idx{2}{二次の部分}{quadratic}{275–6}
\idx{1}{開弦}{open string}{188–9, 196–7, 276}
\idx{1}{弦の場}{string field}{311–14, \textit{327–8}}
\idx{1}{タッドポール項}{tadpole terms}{197–8, 225, 309}
\idx{0}{作用（世界面）}{action (world-sheet)}{}
\idx{1}{$bc$}{}{50}
\idx{1}{$\beta\gamma$}{}{52}
\idx{1}{Brink--Di Vecchia--Howe--Deser--Zumino}{}{12}
\idx{1}{Euclid 作用}{Euclidean}{82–3}
\idx{1}{Euler 項（$\chi$）}{Euler term ($\chi$)}{15–16, \textit{30}, 83–4, 179, 203}
\idx{1}{Faddeev--Popov ゴースト}{Faddeev--Popov ghost}{89, 157–8}
\idx{1}{ゲージ固定}{gauge-fixing}{126–7, 131}
\idx{1}{Hilbert 作用（重力）}{Hilbert (gravitational)}{15–16}
\idx{1}{線形ディラトン}{linear dilaton}{49}
\idx{1}{Liouville}{}{324–5}
\idx{1}{質量ゼロのスカラー}{massless scalar}{32, 34, \textit{346}}
\idx{1}{Nambu--Goto}{}{10–13, 15, 28, \textit{30}}
\idx{2}{Dブレーンの作用との関係}{and D-brane action}{270}
\idx{1}{非臨界弦}{noncritical string}{325–6}
\idx{1}{非線形シグマ模型}{nonlinear sigma model}{108–10, 112}
\idx{2}{$T$双対性}{$T$-duality}{\textit{280–1}}
\idx{1}{Polyakov}{}{12–16, 28}
\idx{2}{非臨界弦との関係}{and noncritical string}{325}
\idx{2}{Euclid 作用}{Euclidean}{82–3}
\idx{2}{平坦な世界面}{flat world-sheet}{32}
\idx{1}{巻き付いたトーラス}{wrapped torus}{237}
\idx{2}{反対称テンソル背景}{antisymmetric tensor background}{249–50}
\idx{0}{アフィン Lie 代数}{affine Lie algebra}{243}
\idx{0}{環帯}{annulus}{円筒（cylinder）を参照}
\idx{0}{異常}{anomaly}{}
\idx{1}{BRST}{}{132–3, \textit{143}}
\idx{2}{弦ループによるもの}{from string loops}{308–9}
\idx{1}{共形異常}{conformal}{Weyl と同じ}
\idx{1}{カレント保存}{current conservation}{\textit{119}, 158–9}
\idx{1}{ゴースト数}{ghost number}{158–9, 177}
\idx{1}{重力（diff）}{gravitational (diff)}{95, 349}
% Original index: p.390
\idx{1}{Lorentz}{}{23–4, 94}
\idx{1}{トレース異常}{trace}{Weyl と同じ}
\idx{1}{Weyl}{}{90–7, \textit{119}, 349}
\idx{2}{弦ループによるもの}{from string loops}{308–9}
\idx{2}{非線形シグマ模型におけるもの}{in nonlinear sigma model}{111–12}
\idx{2}{頂点演算子におけるもの}{in vertex operators}{102–7}
\idx{2}{積分したもの}{integrated}{93–4}
\idx{0}{反交換する c 数}{anticommuting c-numbers}{341}
\idx{0}{反周期境界条件}{antiperiodic boundary conditions}{}
\idx{1}{トレースにおけるもの}{in traces}{212, 336, 344, \textit{345–6}}
\idx{1}{ねじれた状態におけるもの}{in twisted states}{257–9}
\idx{0}{反対称テンソル場}{antisymmetric tensor field}{}
\idx{1}{アクシオンとして}{as axion}{29}
\idx{1}{背景}{background}{108–15, \textit{120}}
\idx{2}{定数の背景}{constant}{249–51, 254–5}
\idx{1}{結合先}{coupling to}{}
\idx{2}{Dブレーン}{D-brane}{270–2, \textit{282}}
\idx{2}{巻き付いた弦}{winding strings}{241–2}
\idx{1}{次元縮約}{dimensional reduction}{234–5, 241}
\idx{1}{場の方程式}{field equation}{111–12}
\idx{1}{ゲージ不変性}{gauge invariance}{27, 109, 196}
\idx{1}{BRST 量子化におけるもの}{in BRST quantization}{137}
\idx{1}{光円錐量子化におけるもの}{in light-cone quantization}{27–30}
\idx{1}{OCQ におけるもの}{in OCQ}{125}
\idx{1}{時空の作用}{spacetime action}{114}
\idx{1}{頂点演算子}{vertex operator}{104–7}
\idx{0}{アクシオン}{axion}{1, 6, 29}
\idx{0}{$bc$ CFT}{}{50–1, \textit{75–6}, 89}
\idx{1}{期待値}{expectation values}{176–8, 212–13}
\idx{1}{モード展開}{mode expansion}{61–3}
\idx{1}{状態・演算子写像}{state--operator mapping}{65–6}
\idx{0}{$\beta\gamma$ CFT}{}{52, \textit{75}}
\idx{0}{BRST 不変性}{BRST invariance}{}
\idx{1}{弦の振幅の不変性}{of string amplitudes}{160–1, 283, 285, 287–90, 351}
\idx{1}{頂点演算子の不変性}{of vertex operators}{164, \textit{165}}
\idx{0}{BRST 量子化}{BRST quantization}{126ff, 350}
\idx{1}{弦の場の理論との関係}{and string field theory}{311–15}
\idx{1}{冪零性}{nilpotence}{128, 130–2, 138, \textit{143}}
\idx{1}{点粒子}{point particle}{129–31}
\idx{1}{弦}{string}{131–43}
\idx{2}{異常}{anomaly}{132, \textit{143}, 308–9}
\idx{2}{電荷}{charge}{132}
\idx{2}{カレント}{current}{131}
\idx{2}{光円錐量子化との同値性}{equivalence to light-cone}{137–41}
\idx{2}{OCQ との同値性}{equivalence to OCQ}{141–3}
\idx{2}{ノーゴースト定理}{no-ghost theorem}{137–41}
\idx{2}{変換}{transformation}{131}
\idx{0}{背景電荷}{background charge}{170, \textit{204}, 208}
\idx{0}{背景場}{background fields}{108–18}
\idx{1}{線形化と頂点演算子}{linearized, and vertex operators}{108–11}
\idx{0}{Batalin--Vilkovisky（BV）形式}{Batalin--Vilkovisky (BV) formalism}{128}
\idx{0}{Beltrami 微分}{Beltrami differential}{162}
\idx{0}{Beltrami 方程式}{Beltrami's equation}{162}
\idx{0}{Berezin 積分}{Berezin integration}{341}
\idx{1}{Jacobian}{}{345}
\idx{0}{ベータ関数（Euler）}{beta function (Euler)}{183}
\idx{0}{ベータ関数（場の理論）}{beta function (field theory)}{91, 111–12, 349}
\idx{1}{弦ループ補正}{string loop corrections}{308–9}
\idx{1}{世界面補正}{world-sheet corrections}{112, 195}
\idx{0}{ブラックホール}{black holes}{6–7, 320}
\idx{0}{Boltzmann 重み}{Boltzmann weight}{325}
\idx{0}{Borel 総和}{Borel summation}{316, 352}
\idx{0}{Born--Infeld 作用}{Born--Infeld action}{271, 273}
\idx{0}{ボソン弦理論}{bosonic string theory}{}
\idx{1}{この理論に特有の非現実的性質}{artifacts of}{23, 219, 223–4, 246, 275, 306}
\idx{1}{研究する理由}{reason for studying}{8}
\idx{0}{ボソン化}{bosonization}{303}
\idx{0}{境界状態}{boundary state}{226, \textit{230}}
\idx{0}{$CPT$ 対称性}{$CPT$ symmetry}{23}
\idx{0}{Calabi--Yau コンパクト化}{Calabi--Yau compactification}{352}
\idx{0}{Campbell--Baker--Hausdorff 公式}{Campbell--Baker--Hausdorff formula}{202}
\idx{0}{伝播関数の相殺の議論}{canceled propagator argument}{161, 196, 288, 309}
\idx{0}{Casimir エネルギー}{Casimir energy}{22}
\idx{1}{ユニタリ CFT におけるもの}{in unitary CFT}{73}
\idx{1}{覚え方}{mnemonic}{73}
\idx{0}{中心電荷}{central charge}{48, \textit{75–6}}
\idx{1}{BRST 異常との関係}{and BRST anomaly}{132, \textit{143}}
\idx{1}{コンパクト化との関係}{and compactification}{113, 117, 137}
\idx{1}{状態密度との関係}{and density of states}{213}
\idx{1}{重力異常との関係}{and gravitational anomaly}{95}
\idx{1}{基底状態エネルギーとの関係}{and ground state energy}{73}
\idx{1}{Weyl 異常との関係}{and Weyl anomaly}{92–6}
% Original index: p.391
\idx{1}{$bc$}{}{51}
\idx{1}{$\beta\gamma$}{}{52}
\idx{1}{自由スカラー}{free scalar}{48}
\idx{1}{Virasoro 代数におけるもの}{in Virasoro algebra}{56}
\idx{1}{線形ディラトン}{linear dilaton}{49}
\idx{1}{正値性}{positivity}{72, \textit{76}}
\idx{0}{Chan--Paton 因子}{Chan--Paton factor}{186–7}
\idx{0}{Chan--Paton 状態}{Chan--Paton states}{30, 184–92}
\idx{1}{Dブレーンとの関係}{and D-branes}{268–70}
\idx{1}{円筒の振幅におけるもの}{in cylinder amplitude}{223}
\idx{1}{Möbius の帯におけるもの}{in Möbius strip}{228}
\idx{1}{非 Abel ゲージ相互作用}{non-Abelian gauge interactions}{188–9}
\idx{1}{向き付けのない弦}{unoriented string}{189–92}
\idx{0}{Chern--Simons 項}{Chern--Simons term}{234–5}
\idx{0}{カイラルなゲージ対称性}{chiral gauge symmetry}{6}
\idx{0}{閉状態（BRST）}{closed state (BRST)}{128}
\idx{0}{閉弦}{closed string}{14}
\idx{1}{BRST 量子化}{BRST quantization}{136–7}
\idx{1}{光円錐量子化}{light-cone quantization}{25–9}
\idx{1}{OCQ}{}{124–5}
\idx{1}{1ループ振幅}{one-loop amplitudes}{216–22, 226–8}
\idx{1}{円筒の振幅における極}{poles in cylinder amplitude}{225–6, \textit{230}, 274–6}
\idx{1}{ツリーレベル振幅}{tree-level amplitudes}{192–8}
\idx{2}{開弦との関係}{relation to open}{195}
\idx{1}{頂点演算子}{vertex operators}{63–8, 101–7}
\idx{0}{コサイクル}{cocycle}{239}
\idx{0}{係数関数（OPE）}{coefficient functions (OPE)}{37–8}
\idx{1}{共形ブートストラップとの関係}{and conformal bootstrap}{304}
\idx{1}{共形不変性との関係}{and conformal invariance}{47}
\idx{1}{行列要素との関係}{and matrix elements}{200}
\idx{1}{3点振幅との関係}{and three-point amplitudes}{284}
\idx{1}{CFT の決定}{determine CFT}{303}
\idx{1}{一般化したもの}{generalized}{303}
\idx{1}{交換対称性}{interchange symmetry}{40–1}
\idx{0}{コホモロジー（BRST）}{cohomology (BRST)}{128–9, 134–43, \textit{144}, 350}
\idx{1}{振幅との関係}{and amplitudes}{283, 285, 289–90}
\idx{0}{Coleman--Weinberg 公式}{Coleman--Weinberg formula}{221, \textit{229–30}, 350}
\idx{0}{集団座標（Dブレーン）}{collective coordinates (D-brane)}{269, \textit{281–2}}
\idx{0}{交換子と反交換子}{commutators and anticommutators}{}
\idx{1}{正準交換関係}{canonical}{17, 20, 26, 130}
\idx{1}{Jacobi 恒等式}{Jacobi identity}{128, 133, \textit{143}}
\idx{1}{OPE との関係}{relation to OPE}{54–5}
\idx{1}{符号が逆のもの}{wrong-sign}{27–8, 72, 121}
\idx{0}{交換子と反交換子（具体例）}{commutators and anticommutators (specific)}{}
\idx{1}{$[\alpha_m^\mu,\alpha_n^\nu]$}{}{20, 26–7, 58, 62}
\idx{1}{$\{b_n,c_n\}$}{}{61, 130, 133}
\idx{1}{$[\delta_{\tau_1},\delta_{\tau_2}]$}{}{129}
\idx{1}{$[G_I,G_J]$（制約）}{$[G_I,G_J]$ (constraints)}{133}
\idx{1}{$[j_m^i,j_n^j]$}{}{243}
\idx{1}{$[L_m,L_n]$}{}{56–7, \textit{75–6}}
\idx{1}{$[L_m,\mathcal O_n]$}{}{56}
\idx{1}{$[N^{\mathrm g},b_m]$}{}{62}
\idx{1}{$[N^{\mathrm g},c_m]$}{}{62}
\idx{1}{$[p_\mu,x^\nu]$}{}{17, 20, 26, 58, 62, 130}
\idx{1}{$[Q,\mathcal A]$}{}{55}
\idx{1}{$[Q_1,Q_2]$}{}{55}
\idx{1}{$\{Q_B,b_m\}$}{}{132}
\idx{1}{$\{Q_B,L_m\}$}{}{\textit{143}}
\idx{1}{$\{Q_B,Q_B\}$}{}{132–3, \textit{143}}
\idx{1}{$[X_L(z),X_L(z')]$}{}{239}
\idx{0}{コンパクトな CFT}{compact CFT}{117, 137}
\idx{0}{コンパクト化}{compactification}{2, 6, 115–18}
\idx{1}{抽象的 CFT}{abstract CFT}{117–18, 137}
\idx{1}{スケール}{scale}{110, 112–13, 116, 244–5}
\idx{1}{}{}{オービフォールド（orbifold）、オリエンティフォールド（orientifold）、トーラス（toroidal）も参照}
\idx{0}{複素座標}{complex coordinates}{33–4, 153}
\idx{1}{発散定理}{divergence theorem}{33–4}
\idx{0}{複素多様体}{complex manifold}{153}
\idx{0}{共形ブートストラップ}{conformal bootstrap}{304}
\idx{0}{共形場理論（CFT）}{conformal field theory (CFT)}{43–76, 349, 351}
\idx{1}{コンパクト化との関係}{and compactification}{115–18}
\idx{1}{$c=1$}{}{261–3, 351}
\idx{1}{コンパクトなもの}{compact}{117, 137}
\idx{1}{曲がった計量への結合}{coupling to curved metric}{95, 97}
\idx{1}{期待値}{expectation values}{}
\idx{2}{行列要素との関係}{and matrix elements}{199–201}
\idx{2}{Möbius 不変性}{Möbius invariance}{198–9, \textit{205}}
\idx{2}{正の種数}{positive genus}{169–78, 201–2}
\idx{2}{種数ゼロ}{zero genus}{208–16}
\idx{1}{自由なもの}{free}{32–6, 49–52}
\idx{1}{物質}{matter}{122, 131}
\idx{1}{モード展開}{mode expansions}{58–63}
\idx{1}{Riemann面上のもの}{on Riemann surface}{154}
\idx{1}{周期的スカラー}{periodic scalar}{235–40, 242–4, 247–8}
% Original index: p.392
\idx{1}{縫い合わせと切断}{sewing and cutting}{300–5}
\idx{1}{状態・演算子同型}{state--operator isomorphism}{63–8}
\idx{1}{ユニタリなもの}{unitary}{71–3, 117, 137}
\idx{1}{}{}{$bc$、$\beta\gamma$、自由スカラー（free scalar）、線形ディラトン（linear dilaton）の各 CFT も参照}
\idx{0}{共形ゲージ}{conformal gauge}{85}
\idx{1}{光円錐ゲージとの関係}{and light-cone gauge}{83}
\idx{1}{曲率}{curvature}{93}
\idx{1}{ゴースト作用}{ghost action}{89}
\idx{1}{OCQ}{}{121}
\idx{1}{球面}{sphere}{166}
\idx{1}{周囲の変分}{variations around}{152}
\idx{0}{共形不変性}{conformal invariance}{43–5}
\idx{1}{OPE との関係}{and OPE}{45–7}
\idx{1}{diff$\times$Weyl 不変性との関係}{relation to diff$\times$Weyl invariance}{85–6}
\idx{0}{共形 Killing 方程式}{conformal Killing equation}{151}
\idx{0}{共形 Killing 群（CKG）またはベクトル（CKV）}{conformal Killing group (CKG) or vector (CKV)}{149–52, 155, 157, 161}
\idx{1}{具体的な Riemann面}{specific Riemann surfaces}{}
\idx{2}{正の種数}{positive genus}{167–9, 176, 178, 192, 197}
\idx{2}{種数ゼロ}{zero genus}{206–8, 212, 216–17}
\idx{0}{共形変換}{conformal transformation}{44–5, \textit{75}}
\idx{1}{Virasoro 代数との関係}{and Virasoro algebra}{52–8}
\idx{1}{$d\geq3$}{}{\textit{75}}
\idx{1}{エネルギー運動量テンソルの変換}{of energy-momentum tensor}{47–9}
\idx{1}{ゴースト数カレントの変換}{of ghost-number current}{51}
\idx{0}{共役類}{conjugacy classes}{260}
\idx{0}{制約（量子化におけるもの）}{constraints (in quantization)}{28, 130, 133, \textit{143–4}, 350}
\idx{0}{縮約}{contraction}{39–40, 50, 65, \textit{75}, 172, 175}
\idx{1}{自己縮約}{self-}{103, 171, 338}
\idx{0}{座標変換}{coordinate transformation}{}
\idx{1}{大きな変換}{large}{モジュラー不変性（modular invariance）を参照}
\idx{1}{時空}{spacetime}{28, 109, 196, 285, 310}
\idx{1}{世界面}{world-sheet}{diff 不変性（diff invariance）を参照}
\idx{0}{宇宙定数}{cosmological constant}{}
\idx{1}{時空}{spacetime}{222, 306, 350}
\idx{1}{世界面}{world-sheet}{95, 323}
\idx{0}{共変微分}{covariant derivative}{89, \textit{118}}
\idx{0}{共変ゲージ}{covariant gauge}{共形ゲージ（conformal gauge）を参照}
\idx{0}{臨界次元}{critical dimension}{}
\idx{1}{背景の効果}{effect of backgrounds}{113, \textit{120}}
\idx{1}{BRST 異常から}{from BRST anomaly}{132, \textit{143}}
\idx{1}{円筒の振幅から}{from cylinder amplitude}{\textit{230}}
\idx{1}{Lorentz 異常から}{from Lorentz anomaly}{21–4, 27}
\idx{1}{モジュラー不変性から}{from modular invariance}{\textit{230}}
\idx{1}{OCQ から}{from OCQ}{124, \textit{143}}
\idx{1}{Veneziano 振幅から}{from Veneziano amplitude}{\textit{204}}
\idx{1}{Weyl 異常から}{from Weyl anomaly}{94–5}
\idx{0}{臨界点}{critical point}{44}
\idx{0}{複比}{cross-ratio}{199}
\idx{0}{クロスキャップ}{cross-cap}{101–2, 207–8, 227–8}
\idx{1}{状態}{state}{\textit{230}}
\idx{0}{カレント代数}{current algebra}{243–4}
\idx{0}{曲がった時空}{curved spacetime}{108–18}
\idx{0}{サイクル}{cycle}{}
\idx{1}{$A$ および $B$ サイクル}{$A$- and $B$-}{147–8, 292, 296, 304}
\idx{0}{円筒（環帯）}{cylinder (annulus)}{78, 81}
\idx{1}{長いもの}{long}{ハンドル、退化（handle, degenerating）を参照}
\idx{1}{Riemann面}{Riemann surface}{207}
\idx{1}{半無限のもの}{semi-infinite}{63–4}
\idx{1}{弦の振幅}{string amplitude}{222–6, \textit{229–30}}
\idx{2}{Dブレーン張力の計算}{calculation of D-brane tension}{274–6}
\idx{2}{閉弦の極}{closed string poles}{225–6, \textit{230}}
\idx{2}{発散}{divergence}{224–5, 229}
\idx{2}{向き付けのないもの}{unoriented}{226}
\idx{0}{Dブレーン}{D-brane}{7, 31, 104, 268–82, 317, 351}
\idx{1}{作用}{action}{270–3}
\idx{1}{ゲージ対称性}{gauge symmetry}{268–73}
\idx{1}{半Dブレーン}{half-D-brane}{280}
\idx{1}{張力}{tension}{273–7}
\idx{0}{Dedekind イータ関数}{Dedekind eta function}{210, 216}
\idx{0}{Del Guidice--Di Vecchia--Fubini（DDF）演算子}{Del Guidice--Di Vecchia--Fubini (DDF) operators}{143, 240}
\idx{0}{状態密度}{density of states}{213, 321}
\idx{0}{行列式}{determinant}{}
\idx{1}{汎関数行列式}{functional}{182, 338–41, 344–5}
\idx{1}{Born--Infeld 作用におけるもの}{in Born--Infeld action}{271}
\idx{1}{Laplacian の行列式}{of Laplacian}{170}
\idx{1}{変分}{variation}{12}
\idx{1}{}{}{Faddeev--Popov 行列式（Faddeev--Popov determinant）も参照}
\idx{0}{diff（再パラメータ化）不変性}{diff (reparameterization) invariance}{12–13}
\idx{1}{異常}{anomaly}{91, 95, 349}
\idx{1}{ゲージ固定}{gauge-fixing}{85, 129, 325}
% Original index: p.393
\idx{1}{頂点演算子の不変性}{of vertex operators}{102–3}
\idx{1}{共形不変性との関係}{relation to conformal invariance}{44, 85–6}
\idx{1}{非物理的状態の除去}{removes unphysical states}{27–8, 117}
\idx{1}{境界がある場合}{with boundary}{90}
\idx{0}{ディラトン}{dilaton}{6}
\idx{1}{等価原理との関係}{and equivalence principle}{115}
\idx{1}{弦の結合定数との関係}{and string coupling}{112–15}
\idx{1}{背景}{background}{108–15, 233–4}
\idx{1}{結合先}{coupling to}{}
\idx{2}{Dブレーン}{D-brane}{270–7}
\idx{2}{開弦}{open strings}{196–8}
\idx{2}{オリエンティフォールド平面}{orientifold plane}{280}
\idx{1}{次元縮約}{dimensional reduction}{233–4, 241}
\idx{1}{$T$双対性の効果}{effect of $T$-duality on}{248}
\idx{1}{場の方程式}{field equation}{111–12, 246}
\idx{1}{BRST 量子化におけるもの}{in BRST quantization}{137}
\idx{1}{光円錐量子化におけるもの}{in light-cone quantization}{27, 29–30}
\idx{1}{OCQ におけるもの}{in OCQ}{125}
\idx{1}{線形ディラトン}{linear}{線形ディラトン CFT（linear dilaton CFT）を参照}
\idx{1}{モジュライ空間}{moduli space}{254}
\idx{1}{伝播関数（場の理論）}{propagator (field theory)}{276}
\idx{1}{時空の作用}{spacetime action}{114–15, 196–8, 225, 275–6}
\idx{1}{頂点演算子}{vertex operator}{104–7}
\idx{0}{次元縮約}{dimensional reduction}{235, 249}
\idx{0}{Dirichlet 境界条件}{Dirichlet boundary condition}{\textit{31}, 266–70, \textit{281}}
\idx{1}{Neumann 条件との $T$ 双対性}{$T$-dual to Neumann}{266–70}
\idx{0}{離散ゲージ対称性}{discrete gauge symmetry}{260–1, 277–8}
\idx{0}{円板}{disk}{}
\idx{1}{半無限円筒との共形同値性}{conformal to semi-infinite cylinder}{63–8}
\idx{1}{期待値}{expectation values}{174–5, 178}
\idx{1}{Riemann面}{Riemann surface}{168}
\idx{1}{弦の振幅}{string amplitudes}{178–89, 196–8, \textit{204–5}}
\idx{2}{Dブレーンを伴うもの}{with D-branes}{270, 274, \textit{281–2}}
\idx{1}{タッドポール}{tadpole}{197–8, 223–6}
\idx{2}{射影平面との相殺}{cancellation with projective plane}{229}
\idx{0}{発散}{divergences}{}
\idx{1}{発散を越える解析接続}{analytic continuation past}{181, 287}
\idx{1}{$SO(8192)$ での相殺}{cancellation for $SO(8192)$}{228–9}
\idx{1}{モジュライ空間の境界に由来するもの}{from boundary of moduli space}{196, 296, 305–6}
\idx{1}{IR}{}{181, 223–5, 306–10}
\idx{1}{UV}{}{2–7, 218–20}
\idx{2}{弦理論での不在}{absence in string theory}{4–5, 79, 219, 223–5, 309–10}
\idx{1}{ゼロモード}{zero mode}{173}
\idx{0}{倍化の手法}{doubling trick}{57, 63, 66, 178}
\idx{0}{双対共鳴模型}{dual resonance model}{348}
\idx{0}{双対性}{duality}{$T$双対性（$T$-duality）、世界面双対性（world-sheet duality）を参照}
\idx{0}{有効場の理論}{effective field theory}{110, 114, 117, 188–9, 196–7, 270–3, 309}
\idx{0}{Einstein 計量}{Einstein metric}{115, 194, 275}
\idx{0}{Einstein 方程式}{Einstein's equation}{112–13, 222}
\idx{0}{電弱スケール}{electroweak scale}{222}
\idx{0}{エネルギー運動量テンソル}{energy-momentum tensor}{13, 91}
\idx{1}{モジュライに関する微分との関係}{and derivative with respect to moduli}{161, 210–11}
\idx{1}{制約として}{as constraint}{121}
\idx{1}{境界条件}{boundary condition}{56–7}
\idx{1}{選び方}{choice of}{49, 95, 303}
\idx{1}{共形変換}{conformal transformation}{47–9}
\idx{1}{トレース異常}{trace anomaly}{92–7}
\idx{1}{トレースレス性と正則性}{traceless and holomorphic}{43–4}
\idx{1}{Chan--Paton 状態に対する消失}{vanishes for Chan--Paton states}{185}
\idx{0}{拡大されたゲージ対称性}{enhanced gauge symmetry}{242–4, 248, 254}
\idx{0}{イプシロン処方}{epsilon prescription}{182, 287–8, 305–6}
\idx{0}{同値類}{equivalence classes}{}
\idx{1}{BRST}{}{128}
\idx{1}{ゲージ}{gauge}{84, 87}
\idx{1}{OCQ}{}{123}
\idx{0}{等価原理}{equivalence principle}{115}
\idx{0}{Euclid 随伴}{Euclidean adjoint}{202–3, 336}
\idx{0}{Euclid への接続}{Euclidean continuation}{82–3, 170, 210, 335}
\idx{0}{Euler 数}{Euler number}{96, 101, \textit{118}, 152}
\idx{1}{弦の作用におけるもの}{in string action}{83–4, 112, 179, 203}
\idx{0}{完全状態（BRST）}{exact state (BRST)}{128}
\idx{0}{例外 Lie 代数}{exceptional Lie algebras}{192}
\idx{0}{期待値}{expectation values}{34, 36}
\idx{1}{計算}{calculations}{169–78, 201–2, 208–13}
\idx{0}{指数演算子}{exponential operator}{40, 46}
\idx{0}{外微分}{exterior derivative}{128}
% Original index: p.394
\idx{0}{Faddeev--Popov 行列式}{Faddeev--Popov determinant}{86–90, 154–8}
\idx{1}{円板上}{on disk}{176}
\idx{1}{球面上}{on sphere}{176}
\idx{1}{トーラス上}{on torus}{212–13, 216–17}
\idx{1}{簡便法}{shortcuts}{179, 218}
\idx{1}{符号}{sign}{156, 177, 345}
\idx{0}{Faddeev--Popov ゴースト}{Faddeev--Popov ghosts}{51–2, 88–90, 126}
\idx{1}{境界条件}{boundary condition}{90}
\idx{1}{中心電荷}{central charge}{94}
\idx{1}{真空}{vacuum}{61, 141}
\idx{0}{Faddeev--Popov の手続き、一般}{Faddeev--Popov procedure, general}{126}
\idx{0}{Feynman 図}{Feynman diagrams}{}
\idx{1}{Dブレーン間相互作用の計算}{calculation of D-brane interaction}{275–6}
\idx{1}{Weyl 異常の計算}{calculation of Weyl anomaly}{\textit{119}}
\idx{1}{弦の振幅の分解}{decomposition of string amplitudes}{297–306}
\idx{1}{弦の図との比較}{vs string graphs}{3–5, 101}
\idx{0}{基準計量}{fiducial metric}{85}
\idx{0}{場の強さ、三つの添字を持つもの}{field strength, three-index}{109, 112, 234–5}
\idx{0}{第一量子化}{first-quantization}{310}
\idx{0}{Fischler--Susskind 機構}{Fischler--Susskind mechanism}{308–9}
\idx{0}{固定点}{fixed point}{}
\idx{1}{オービフォールド}{orbifold}{256–8}
\idx{1}{UV}{}{4}
\idx{0}{平坦方向}{flat direction}{246, 262–3, 272–3}
\idx{0}{自由エネルギー}{free energy}{321–2, \textit{328}}
\idx{0}{自由パラメータの不在}{free parameters, absence of}{6, 83–4, 112–13}
\idx{0}{自由スカラー CFT}{free scalar CFT}{32–41}
\idx{1}{期待値}{expectation values}{169–76, 208–11, \textit{229}}
\idx{1}{左向き・右向き}{left-, right-moving}{238}
\idx{1}{モード展開}{mode expansion}{58–61}
\idx{1}{状態・演算子写像}{state--operator mapping}{63–5}
\idx{0}{Fuchs 群}{Fuchsian group}{292–3}
\idx{0}{汎関数積分}{functional integral}{経路積分（path integral）を参照}
\idx{0}{基本領域}{fundamental region}{148–50}
\idx{1}{トーラス振幅におけるもの}{in torus amplitude}{216, 219–20}
\idx{0}{ゲージボソン}{gauge boson}{1, 4}
\idx{1}{開弦}{open string}{28, 123–4, 135–6}
\idx{2}{振幅}{amplitudes}{187–9, \textit{205}}
\idx{2}{BRST 量子化におけるもの}{in BRST quantization}{135–6}
\idx{2}{光円錐量子化におけるもの}{in light-cone quantization}{23}
\idx{2}{OCQ におけるもの}{in OCQ}{123–4}
\idx{2}{頂点演算子}{vertex operator}{107}
\idx{1}{閉弦}{closed string}{241–3}
\idx{2}{振幅}{amplitudes}{242}
\idx{0}{ゲージ結合定数}{gauge coupling}{115, 234, 242–5, 276}
\idx{0}{ゲージ固定}{gauge-fixing}{}
\idx{1}{ゲージの選択への依存性}{dependence on gauge choice}{90–1, 121–2, 126–8}
\idx{1}{diff$\times$Weyl}{}{17–18, 25, 84–6, 131, 154–5}
\idx{1}{一般}{general}{126–8}
\idx{1}{}{}{個々のゲージ（specific gauges）も参照}
\idx{0}{ゲージ相互作用}{gauge interaction}{2, 28, 78, 184–92}
\idx{0}{ゲージ変換}{gauge transformation}{107}
\idx{0}{Gaussian 平均をとったゲージ}{Gaussian-averaged gauges}{127}
\idx{0}{一般相対性理論}{general relativity}{1–2, 5, 28, 108}
\idx{1}{オフシェル振幅}{off-shell amplitudes}{104}
\idx{0}{種数}{genus}{100–1, 150}
\idx{1}{大きい種数}{large}{315–17}
\idx{0}{測地的曲率}{geodesic curvature}{\textit{30}, 82, 96}
\idx{0}{測地距離}{geodesic distance}{102}
\idx{0}{測地線方程式}{geodesic equation}{108}
\idx{0}{ゴースト数}{ghost number}{51, 62, 127}
\idx{1}{異常}{anomaly}{\textit{119}, 158–9, 177}
\idx{1}{弦の場のゴースト数}{of string field}{313–14}
\idx{1}{状態と演算子の比較}{states vs operators}{65–6}
\idx{0}{ゴースト}{ghosts}{Faddeev--Popov ゴースト（Faddeev--Popov ghosts）を参照}
\idx{0}{大域対称性}{global symmetry}{189}
\idx{0}{Goldstone ボソン}{Goldstone boson}{233}
\idx{0}{次数付き Lie 代数}{graded Lie algebra}{\textit{143}}
\idx{0}{大統一}{grand unification}{2, 5}
\idx{0}{Grassmann 変数}{Grassmann variables}{35, 88, 156–8, 341–5}
\idx{0}{重力結合定数}{gravitational coupling}{114}
\idx{1}{関係する結合定数}{relation to}{}
\idx{2}{ゲージ結合定数}{gauge coupling}{234, 244–5}
\idx{2}{弦の結合定数}{string coupling}{194, \textit{205}}
\idx{0}{重力子}{graviton}{}
\idx{1}{結合先}{coupling to}{}
\idx{2}{Dブレーン}{D-brane}{270–7}
\idx{2}{開弦}{open strings}{196–8}
\idx{2}{オリエンティフォールド平面}{orientifold plane}{280}
\idx{1}{BRST 量子化におけるもの}{in BRST quantization}{137}
% Original index: p.395
\idx{1}{光円錐量子化におけるもの}{in light-cone quantization}{27, 29–30}
\idx{1}{OCQ におけるもの}{in OCQ}{125}
\idx{1}{伝播関数（場の理論）}{propagator (field theory)}{276}
\idx{1}{頂点演算子}{vertex operator}{104–7}
\idx{0}{重力}{gravity}{1–7, 78}
\idx{0}{Green 関数}{Green's function}{169–71, 174–6, 208, \textit{229–30}, 275}
\idx{0}{Hagedorn 温度}{Hagedorn temperature}{321–2, 352}
\idx{0}{Hamiltonian}{}{}
\idx{1}{経路積分との関係}{and path integral}{329–32, 335–6, 343}
\idx{1}{背景の効果}{effect of backgrounds}{250, 264}
\idx{1}{分配関数におけるもの}{in partition function}{209, 212–13}
\idx{1}{点粒子}{point particle}{17, 130, 218}
\idx{1}{弦}{string}{19–21, 25–6}
\idx{2}{$L_0$ との関係}{relation to $L_0$}{54, 73, 125}
\idx{0}{ハンドル}{handle}{83, 100–1, 290–2}
\idx{1}{退化}{degenerating}{97–8, 101, 224–6, 229, 295–6, 305–9}
\idx{0}{硬い散乱}{hard scattering}{183–4, 194, \textit{281–2}, 317–19, 352}
\idx{0}{ヘテロティック弦}{heterotic string}{348}
\idx{0}{隠れたセクター}{hidden sector}{6}
\idx{0}{Higgs 場}{Higgs field}{1, 222, 245–6}
\idx{0}{高エネルギー散乱}{high-energy scattering}{硬い散乱（hard scattering）、Regge 散乱（Regge scattering）を参照}
\idx{0}{高温}{high temperature}{320–2}
\idx{0}{最高ウェイト状態}{highest weight state}{70–2, 122, 125, \textit{143}}
\idx{1}{テンソル（プライマリー）演算子との関係}{and tensor (primary) operators}{71}
\idx{0}{正則}{holomorphic}{}
\idx{1}{カレント}{current}{42}
\idx{1}{場}{field}{34, 56, 72}
\idx{1}{二次微分}{quadratic differential}{152}
\idx{1}{遷移関数}{transition functions}{153–4}
\idx{1}{ベクトル場}{vector field}{152}
\idx{0}{超楕円曲面}{hyperelliptic surface}{294}
\idx{0}{赤外（IR）発散}{infrared (IR) divergences}{181, 223–5, 306–10}
\idx{0}{継承原理}{inheritance principle}{258, 261}
\idx{0}{内積}{inner product}{}
\idx{1}{双線形内積（Zamolodchikov）}{bilinear (Zamolodchikov)}{199–204}
\idx{2}{縫い合わせにおけるもの}{in sewing}{301–2}
\idx{2}{Hermite 内積との関係}{relation to Hermitean}{202–4}
\idx{1}{Fermi 振動子}{Fermi oscillator}{343}
\idx{1}{Hermite 内積}{Hermitean}{134–5}
\idx{1}{正の内積}{positive}{71–2, 116–17, 138}
\idx{1}{トレースレス対称テンソル}{traceless symmetric tensor}{\textit{118}}
\idx{0}{インスタントン}{instantons}{316}
\idx{0}{相互作用}{interaction}{}
\idx{1}{非重要な相互作用}{irrelevant}{3}
\idx{1}{弦}{string}{77–82, 104, 178ff}
\idx{0}{Jacobi 恒等式}{Jacobi identity}{128, 133, \textit{143}}
\idx{0}{Jacobi の難解な恒等式}{Jacobi's abstruse identity}{215}
\idx{0}{Jacobian、Berezin 積分}{Jacobian, Berezin integral}{345}
\idx{0}{Kac--Moody 代数}{Kac--Moody algebra}{243}
\idx{0}{Kaluza--Klein ゲージ場}{Kaluza--Klein gauge field}{2, 232–4, 241–3, 249, 257, 351}
\idx{0}{Kaluza--Klein 状態}{Kaluza--Klein states}{232, 245}
\idx{0}{Klein の壺}{Klein bottle}{}
\idx{1}{Riemann面}{Riemann surface}{207}
\idx{1}{弦の振幅}{string amplitude}{226–9, \textit{230}, 278, 280}
\idx{0}{高次での挙動}{large order behavior}{315–17}
\idx{0}{格子}{lattice}{}
\idx{1}{偶自己双対格子}{even self-dual}{252–3, \textit{281}}
\idx{1}{運動量格子}{momentum}{251–3, \textit{281}}
\idx{0}{Laurent 展開}{Laurent expansion}{モード展開（mode expansion）を参照}
\idx{0}{長さ、最小のもの}{length, minimum}{248, 318}
\idx{0}{レベル}{level}{}
\idx{1}{カレント代数}{current algebra}{243–4}
\idx{1}{振動子}{oscillator}{23, 26–7}
\idx{0}{レベル整合}{level-matching}{27, 258}
\idx{0}{光円錐座標}{light-cone coordinates}{16}
\idx{0}{光円錐ゲージ}{light-cone gauge}{17}
\idx{1}{オフシェル}{off-shell}{104}
\idx{0}{光円錐量子化}{light-cone quantization}{}
\idx{1}{点粒子}{point particle}{17}
\idx{1}{弦}{string}{17–30, 349}
\idx{2}{BRST および OCQ との同値性}{equivalence to BRST and OCQ}{137–43}
\idx{0}{線形ディラトン理論}{linear dilaton theory}{49–50, 61, \textit{75}, \textit{204}, 213, 349}
\idx{1}{背景として}{as background}{113, \textit{120}, 323–4, \textit{328}}
\idx{0}{Liouville 場の理論}{Liouville field theory}{325}
\idx{0}{ループ展開}{loop expansion}{115}
\idx{0}{Lorentz 不変性}{Lorentz invariance}{12, 106–7, 116–17, 290}
\idx{1}{光円錐ゲージにおけるもの}{in light-cone gauge}{20–4, 90, 94, 349}
% Original index: p.396
\idx{0}{Mandelstam 変数}{Mandelstam variables}{180}
\idx{0}{多様体}{manifold}{}
\idx{1}{複素多様体}{complex}{153}
\idx{1}{可微分多様体}{differentiable}{153}
\idx{1}{Riemann 多様体}{Riemannian}{153}
\idx{0}{質量殻条件}{mass-shell condition}{}
\idx{1}{頂点演算子から}{from vertex operators}{101, 106}
\idx{1}{BRST におけるもの}{in BRST}{131, 134, 137–8}
\idx{1}{光円錐量子化におけるもの}{in light-cone quantization}{17, 22, 26}
\idx{0}{行列要素、経路積分との関係}{matrix elements, relation to path integral}{199–201, 332–4}
\idx{0}{行列模型}{matrix models}{326–7, 352}
\idx{0}{モジュライ空間上の測度}{measure on moduli space}{154–65}
\idx{1}{BRST 不変性}{BRST invariance}{160–1}
\idx{1}{導出}{derivation}{154–7}
\idx{1}{ゲージ不変性}{gauge invariance}{159–60}
\idx{2}{表示に用いるもの}{in terms of}{}
\idx{3}{行列式}{determinants}{157–8}
\idx{3}{遷移関数}{transition functions}{161–4}
\idx{1}{}{}{Faddeev--Popov も参照}
\idx{0}{計量（時空）}{metric (spacetime)}{9, 108–18}
\idx{1}{次元縮約}{dimensional reduction}{231–3}
\idx{1}{Einstein}{}{115, 194, 275}
\idx{1}{内部計量}{internal}{249ff}
\idx{1}{光円錐座標}{light-cone coordinates}{16}
\idx{1}{シグマ模型}{sigma model}{115}
\idx{1}{弦}{string}{115}
\idx{0}{計量（世界線）}{metric (world-line)}{10}
\idx{1}{ゲージ固定}{gauge-fixing}{129, 145–6}
\idx{0}{計量（世界面）}{metric (world-sheet)}{}
\idx{1}{一定曲率}{constant curvature}{292}
\idx{1}{平坦な計量}{flat}{32–3}
\idx{1}{ゲージ固定}{gauge-fixing}{17, 84–6, 146–52}
\idx{1}{誘導計量}{induced}{11, 15}
\idx{1}{極小面積}{minimal area}{300, 351}
\idx{1}{非臨界弦}{noncritical string}{94–5, 325–6}
\idx{1}{Polyakov}{}{12–13ff}
\idx{0}{覚え方、零点のもの}{mnemonics, zero-point}{73–4, 125}
\idx{0}{Möbius 群}{Möbius group}{56, 167}
\idx{0}{Möbius の帯}{Möbius strip}{}
\idx{1}{Riemann面}{Riemann surface}{208}
\idx{1}{弦の振幅}{string amplitude}{228–9, 280}
\idx{0}{Möbius 変換}{Möbius transformation}{178–80, 192, 198–9}
\idx{0}{モード（または Laurent）展開}{mode (or Laurent) expansion}{}
\idx{1}{$bc$ CFT}{}{61}
\idx{1}{自由スカラー（$X^\mu$）}{free scalar ($X^\mu$)}{26, 58}
\idx{2}{周期的なもの}{periodic}{235–8}
\idx{2}{ねじれたもの}{twisted}{257}
\idx{1}{正則演算子}{holomorphic operator}{56}
\idx{1}{開弦}{open strings}{20, 57, 62–3}
\idx{2}{Dirichlet}{}{267}
\idx{1}{$T_{ab}$}{}{53, 70}
\idx{0}{模型、可解なもの}{model, solvable}{323–7}
\idx{0}{モジュラー群}{modular group}{148–9, 151}
\idx{1}{テータ関数との関係}{and theta functions}{214–16}
\idx{1}{時空}{spacetime}{253–5}
\idx{1}{モジュラー群を持たない曲面}{surfaces without}{207–8, 223, 278}
\idx{1}{漸近挙動の解析への使用}{used to analyze asymptotics}{213, 224, 227–8}
\idx{0}{モジュラー不変性}{modular invariance}{210, 217–18, \textit{229–30}}
\idx{1}{それによる制約}{constraints from}{213, 304, 350}
\idx{1}{オービフォールド・コンパクト化におけるもの}{in orbifold compactification}{257, 259}
\idx{1}{トーラス・コンパクト化におけるもの}{in toroidal compactification}{237, 251–3}
\idx{1}{ねじりによる構成におけるもの}{in twisting construction}{260}
\idx{1}{整合性に必要なもの}{needed for consistency}{220}
\idx{1}{1点関数}{one-point function}{304}
\idx{0}{モジュライ（時空）}{moduli (spacetime)}{233–4, 249, 265}
\idx{1}{モジュライ空間}{moduli space}{253–5}
\idx{2}{分枝}{branches}{262}
\idx{2}{$c=1$}{}{261–3, 351}
\idx{2}{その上の計量}{metric on}{254}
\idx{0}{モジュライ（世界面）}{moduli (world-sheet)}{}
\idx{1}{遷移関数による表現}{encoded in transition functions}{152–4, 161–4}
\idx{1}{計量のモジュライ}{for metric}{147–52}
\idx{1}{モジュライ空間}{moduli space}{148–52, 290–4}
\idx{2}{その境界}{boundary of}{161, 295–6, 300；CFT の漸近挙動（CFT asymptotics）302–3；境界からの発散（divergences from）196, 224–6, 300, 305–9}
\idx{2}{図式による分解}{graphical decomposition}{297–300, 351}
\idx{2}{開弦の場の理論におけるもの}{in open string field theory}{313–4, 316–7}
\idx{2}{その上の測度}{measure on}{154–65, 350}
\idx{2}{その上の全微分}{total derivative on}{161, 196, 220, 288}
% Original index: p.397
\idx{1}{頂点演算子の位置}{vertex operator positions}{150–1, 163–4}
\idx{0}{Nambu–Goto 作用}{Nambu–Goto action}{10–13, 15, 28, \textit{30}}
\idx{0}{Narain コンパクト化}{Narain compactification}{251–5, 351}
\idx{0}{自然さの問題}{naturalness problem}{2}
\idx{0}{Neumann 境界条件}{Neumann boundary condition}{14, \textit{31}, 62, 174}
\idx{1}{Dirichlet 条件と T 双対}{T-dual to Dirichlet}{265–70}
\idx{0}{冪零性}{nilpotence}{BRST を参照}
\idx{0}{ノーゴースト定理}{no-ghost theorem}{137–43, 350}
\idx{0}{Noether の定理}{Noether's theorem}{41–3, \textit{74–5}}
\idx{1}{異常を伴う場合}{anomalous}{159}
\idx{1}{局所的な性質として}{as local property}{86}
\idx{0}{非 Abel ゲージ理論}{non-Abelian gauge theory}{188–92, 243}
\idx{1}{弦の場の理論から}{from string field theory}{313}
\idx{1}{D-brane 上のもの}{on D-brane}{272–3}
\idx{0}{非可換幾何学}{noncommutative geometry}{272}
\idx{0}{非臨界弦理論}{noncritical string theory}{325–7}
\idx{0}{非線形シグマ模型}{nonlinear sigma model}{108–10, \textit{120}, 349}
\idx{0}{非摂動効果}{nonperturbative effects}{234, 248, 269, 310, 315–7}
\idx{1}{$\exp(-1/g_{\mathrm c})$}{}{317, 327}
\idx{0}{正規順序}{normal ordering}{}
\idx{1}{境界でのもの}{boundary}{\textit{75}, 175}
\idx{1}{共形正規順序}{conformal}{36–40, 50, 60}
\idx{1}{定数}{constants}{22–3, 27, \textit{76}}
\idx{2}{直接計算}{direct calculation}{60}
\idx{2}{状態・演算子対応から}{from state–operator mapping}{72}
\idx{2}{Virasoro 代数から}{from Virasoro algebra}{59, 61}
\idx{2}{覚え方}{mnemonics}{73, 125}
\idx{1}{生成・消滅演算子によるもの}{creation–annihilation}{60}
\idx{1}{微分同相不変なもの}{diff-invariant}{102–3}
\idx{1}{異なる形の間の関係}{relation between different forms}{59–60, \textit{76}}
\idx{0}{零状態}{null state}{}
\idx{1}{BRST}{}{128, 136, 311}
\idx{1}{非結合}{decoupling}{\textit{119}, 161, 196, 220, 285–95}
\idx{1}{OCQ}{}{122–5, \textit{143}}
\idx{0}{$O(k,k,\mathbb R)$}{}{252–3}
\idx{0}{$O(k,k,\mathbb Z)$}{}{253}
\idx{0}{オフシェル振幅}{off-shell amplitudes}{97, 103–4}
\idx{0}{旧共変量子化（OCQ）}{old covariant quantization (OCQ)}{121–5}
\idx{1}{BRST および光円錐量子化との同値性}{equivalence to BRST and light-cone}{141–3}
\idx{1}{頂点演算子}{vertex operators}{125, 143, 160–1, 164, \textit{165}}
\idx{0}{1 ループ振幅}{one-loop amplitudes}{150, 206–30}
\idx{1}{有限性}{finiteness}{220}
\idx{1}{場の理論との比較}{vs field theory}{218–20}
\idx{0}{開弦理論}{open string theory}{14}
\idx{1}{BRST 量子化}{BRST quantization}{131–6}
\idx{1}{Dirichlet}{}{266–70}
\idx{1}{光円錐量子化}{light-cone quantization}{17–25}
\idx{1}{閉弦を含まなければならないこと}{must include closed strings}{81–2, 225}
\idx{1}{OCQ}{}{122–5}
\idx{1}{1 ループ振幅}{one-loop amplitudes}{222–6, 228–9}
\idx{1}{木レベル振幅}{tree-level amplitudes}{178–92}
\idx{2}{閉弦との関係}{relation to closed}{195}
\idx{1}{頂点演算子}{vertex operators}{66, 107}
\idx{0}{演算子による計算}{operator calculations}{201–2}
\idx{0}{演算子方程式}{operator equation}{35–6, 333–4}
\idx{0}{演算子積展開（OPE）}{operator product expansion (OPE)}{36–41, \textit{74–6}}
\idx{1}{共形不変性との関係}{and conformal invariance}{45–9}
\idx{1}{Ward 恒等式との関係}{and Ward identity}{42–3}
\idx{1}{結合則}{associativity}{70, 201, 304}
\idx{1}{収束}{convergence}{37–8}
\idx{1}{導出}{derivation}{}
\idx{2}{CFT}{}{68–70}
\idx{2}{自由場の理論}{free field theory}{37–8}
\idx{1}{一般化されたもの}{generalized}{302–3}
\idx{1}{交換対称性}{interchange symmetry}{40–1}
\idx{1}{Laurent 展開}{Laurent expansion}{70–1}
\idx{1}{有限の隔たりでのもの}{nonzero separation}{48}
\idx{1}{交換子との関係}{relation to commutator}{54–7}
\idx{1}{$z$, $\bar z$ 依存性}{$z$-, $\bar z$-dependence}{47}
\idx{0}{演算子積展開（具体例）}{operator product expansions (specific)}{}
\idx{1}{$\mathcal A_i\,\mathcal A_j$}{}{47}
\idx{1}{$b\,b$, $b\,c$, $c\,c$}{}{50}
\idx{1}{$\beta\,\beta$, $\beta\,\gamma$, $\gamma\,\gamma$}{}{52}
\idx{1}{DDF 演算子}{DDF operators}{240}
\idx{1}{$\partial X^\mu\,\partial X_\mu$}{}{211}
\idx{1}{$:e^{ik\cdot X}:\,{:e^{ik'\cdot X}:}$}{}{40, 251}
\idx{1}{$:\mathcal F:\,{:\mathcal G:}$}{}{39}
\idx{1}{$j_{\mathrm B}\,b$, $j_{\mathrm B}\,c$}{}{132}
\idx{1}{$j_{\mathrm B}\,\mathcal O^{\mathrm m}$}{}{132}
\idx{1}{$j_{\mathrm B}\,j_{\mathrm B}$}{}{132}
% Original index: p.398
\idx{1}{$j_{\mathrm B}\,\mathcal V$}{}{\textit{165}}
\idx{1}{$j^i\,j^j$}{}{243}
\idx{1}{$j^\mu\,{:e^{ik\cdot X}:}$}{}{43}
\idx{1}{$\mathcal O_i\,\mathcal O_j$}{}{200}
\idx{1}{$T\,\mathcal A$}{}{45–6}
\idx{1}{$T\,e^{ik\cdot X}$}{}{46}
\idx{1}{$T\,j$}{}{51}
\idx{1}{$T\,j_{\mathrm B}$}{}{133}
\idx{1}{$T\,\mathcal O$}{}{46}
\idx{1}{$T\,T$}{}{47–8}
\idx{1}{$T\,\widetilde T$}{}{47–8}
\idx{1}{$\mathcal V\,\mathcal V'$}{}{284, 289, \textit{327}}
\idx{1}{$X_{\mathrm L}\,X_{\mathrm L}$, $X_{\mathrm L}\,X_{\mathrm R}$, $X_{\mathrm R}\,X_{\mathrm R}$}{}{238}
\idx{1}{$X^\mu\,X^\nu$}{}{38}
\idx{0}{オービフォールド}{orbifold}{256–63, 351}
\idx{1}{固定点}{fixed points}{256–8}
\idx{1}{継承原理}{inheritance principle}{258, 261}
\idx{1}{モジュラー不変性}{modular invariance}{259–60}
\idx{1}{ねじれた状態}{twisted states}{\textit{31}, 256–61}
\idx{2}{頂点演算子}{vertex operators}{258, 351}
\idx{1}{ねじり構成}{twisting construction}{259–63}
\idx{0}{向きのある弦理論}{oriented string theory}{29}
\idx{1}{}{}{向きのない弦理論（unoriented string theory）の補集合を参照}
\idx{0}{オリエンティフォールド}{orientifold}{277–80}
\idx{1}{固定平面}{fixed planes}{278–80}
\idx{1}{局所的な物理は向きをもつ}{local physics oriented}{278}
\idx{0}{両巻の概要}{outline of both volumes}{8–9}
\idx{0}{$PSL(2,\mathbb C)$}{}{168, 192, 197–9, \textit{205}}
\idx{0}{$PSL(2,\mathbb R)$}{}{168, 178, 199, \textit{205}, 255, 292–3}
\idx{0}{$PSL(2,\mathbb Z)$}{}{148, 255}
\idx{0}{パリティ}{parity}{}
\idx{1}{時空}{spacetime}{6, 255, 277}
\idx{1}{世界面}{world-sheet}{29, 200, 255, 345}
\idx{2}{Chan–Paton 因子への作用}{action on Chan–Paton factors}{189–92}
\idx{2}{ゲージ化}{gauging}{277}
\idx{2}{1 ループ振幅におけるもの}{in one-loop amplitudes}{207–8, 226–8}
\idx{0}{分配関数}{partition function}{209}
\idx{1}{$bc$}{}{212}
\idx{1}{一般の CFT}{general CFT}{213}
\idx{1}{オービフォールド}{orbifold}{258–61}
\idx{1}{スカラー CFT}{scalar CFT}{209–11}
\idx{1}{熱的なもの}{thermal}{320–2, \textit{328}}
\idx{1}{トーラスコンパクト化}{toroidal compactification}{237, 250–2, \textit{281}}
\idx{0}{経路（汎関数）積分}{path (functional) integrals}{32–6}
\idx{1}{期待値との関係}{and expectation values}{34}
\idx{1}{正準交換子の導出}{canonical commutators from}{54–5, 58–61}
\idx{1}{Euclid 回転}{Euclidean rotation}{82–3, 170, 210, 335}
\idx{1}{評価}{evaluation}{}
\idx{2}{Weyl 異常の積分によるもの}{by integrating Weyl anomaly}{93–4}
\idx{2}{演算子の方法によるもの}{by operator methods}{201–2, 209–10, 212}
\idx{2}{直接計算}{direct}{157–9, 169–72, 176–7, 337–41, 344–5}
\idx{2}{正則性を用いるもの}{using holomorphicity}{172–4, 177, 210–11}
\idx{1}{存在}{existence}{332}
\idx{1}{フェルミオンのもの}{fermionic}{341–5}
\idx{1}{トレースに対するもの}{for traces}{335–6, 344}
\idx{1}{Hamilton 形式}{Hamiltonian}{331}
\idx{1}{導入}{introduction}{329–46}
\idx{1}{Lagrange 形式}{Lagrangian}{331}
\idx{1}{全微分の積分}{of total derivative}{34, 333, 341}
\idx{1}{その中の演算子方程式}{operator equations in}{35, 333}
\idx{1}{Polyakov}{}{82–4}
\idx{2}{ゲージ固定}{gauge-fixing}{84–90, 154–9}
\idx{2}{曲がった時空におけるもの}{in curved spacetime}{108–12}
\idx{2}{その中の Weyl 異常}{Weyl anomaly in}{90–7}
\idx{1}{切断を通じた演算子表示との関係}{relation to operator expressions via cutting}{55, 199–201, 300–5, 332–4}
\idx{1}{弦振幅の表示}{representation of string amplitudes}{77–82, 97–101}
\idx{1}{ねじれたもの}{twisted}{260, 336}
\idx{1}{状態・演算子対応の導出への利用}{used to derive state–operator mapping}{66–8}
\idx{0}{Pauli–Villars 正則化因子}{Pauli–Villars regulator}{91, 340}
\idx{0}{周期行列}{period matrix}{293}
\idx{0}{周期境界条件}{periodic boundary condition}{14, \textit{31}, 52}
\idx{0}{摂動論、弦の}{perturbation theory, string}{5, 101, 283–310}
\idx{0}{物理状態}{physical state}{}
\idx{1}{BRST}{}{127–9}
% Original index: p.399
\idx{2}{$b_0$ 条件}{$b_0$ condition}{131, 134, 136, 289–90}
\idx{2}{頂点演算子}{vertex operators}{164}
\idx{1}{OCQ}{}{121–3}
\idx{2}{$L_0$ 条件}{$L_0$ condition}{125}
\idx{2}{頂点演算子}{vertex operators}{125, 143, 160–1, 164}
\idx{0}{Planck 長}{Planck length}{3, 116}
\idx{0}{Planck 質量}{Planck mass}{3, 245}
\idx{0}{配管接続}{plumbing fixture}{294–305}
\idx{0}{Poincaré 不変性}{Poincaré invariance}{10, 12–16, 91, 185}
\idx{1}{タドポールとの関係}{and tadpoles}{225}
\idx{1}{4 次元のもの}{four-dimensional}{16}
\idx{1}{境界条件のもの}{of boundary conditions}{14–15}
\idx{0}{点粒子}{point particle}{}
\idx{1}{作用}{action}{9–11}
\idx{1}{Wilson 線との関係}{and Wilson line}{264}
\idx{1}{BRST 量子化}{BRST quantization}{129–31}
\idx{1}{曲がった時空におけるもの}{in curved spacetime}{108}
\idx{1}{光円錐量子化}{light-cone quantization}{17}
\idx{1}{経路積分}{path integral}{77, 82–3, \textit{164–5}}
\idx{2}{真空ループ}{vacuum loop}{145–6, \textit{164}, 218, 220–2}
\idx{0}{Poisson 再和公式}{Poisson resummation formula}{237, 250, 252}
\idx{0}{Polyakov 作用}{Polyakov action}{作用（action）、Polyakov を参照}
\idx{0}{Polyakov 経路積分}{Polyakov path integral}{経路積分（path integral）、Polyakov を参照}
\idx{0}{ポテンシャルエネルギー}{potential energy}{}
\idx{1}{$c=1$ のモジュライ}{$c=1$ moduli}{246}
\idx{1}{ディラトンのタドポールによるもの}{from dilaton tadpoles}{197–8, 225, 309}
\idx{1}{4 次の Yang–Mills 相互作用によるもの}{from quartic Yang–Mills interaction}{265, 272}
\idx{0}{プライマリー場}{primary field}{テンソル演算子（tensor operator）を参照}
\idx{0}{射影平面}{projective plane}{}
\idx{1}{期待値}{expectation values}{176, 178}
\idx{1}{Riemann 面}{Riemann surface}{168–9}
\idx{1}{弦振幅}{string amplitudes}{196–8}
\idx{1}{タドポール}{tadpole}{197–8, 227–8}
\idx{2}{円板との相殺}{cancellation with disk}{229}
\idx{0}{伝播関数}{propagator}{}
\idx{1}{ディラトン}{dilaton}{275–6}
\idx{1}{重力子}{graviton}{275–6}
\idx{1}{光子}{photon}{290}
\idx{1}{点粒子}{point particle}{115, \textit{164–5}}
\idx{1}{極}{pole}{225–6, 307–8}
\idx{1}{弦}{string}{297, 305–6, 313–14}
\idx{0}{陽子崩壊}{proton decay}{1}
\idx{0}{穿孔点}{puncture}{296–300}
\idx{0}{QCD 弦}{QCD string}{25}
\idx{0}{量子化}{quantization}{}
\idx{1}{正準量子化}{canonical}{17, 19–20, 26, 130, 249–50, 164}
\idx{1}{経路積分量子化}{path integral}{32–6, 77–82, 329–46}
\idx{2}{}{}{BRST、光円錐（light-cone）、OCQ、経路積分（path integrals）も参照}
\idx{0}{場の量子論}{quantum field theory}{1–7, 104, 113, 310, 331–2}
\idx{1}{紫外発散}{UV divergence}{218–20}
\idx{1}{世界面}{world-sheet}{13–15, 32ff, 108–10, 118}
\idx{0}{量子重力}{quantum gravity}{2–7, 25}
\idx{0}{量子力学}{quantum mechanics}{1, 116, 329–46}
\idx{0}{Regge 振る舞い}{Regge behavior}{183–4, 193, 320, 352}
\idx{0}{Regge 勾配}{Regge slope}{11, 24, 112}
\idx{0}{繰り込み}{renormalization}{38, 102–3, 189, 332, 340}
\idx{0}{繰り込み群方程式}{renormalization group equation}{4, 91, 111}
\idx{0}{Ricci 曲率}{Ricci curvature}{}
\idx{1}{時空}{spacetime}{111–12, 114}
\idx{2}{次元低減}{dimensional reduction}{232–3}
\idx{1}{世界面}{world-sheet}{15, 85, 92–3}
\idx{0}{Riemann–Roch の定理}{Riemann–Roch theorem}{152, 158–9, 293}
\idx{0}{Riemann 面}{Riemann surface}{100–1, 152–4, 161–4, 351}
\idx{1}{縫い合わせ}{sewing}{294–300, 351}
\idx{1}{具体例}{specific}{166–9, 206–8, 290–4}
\idx{0}{Riemann テンソル}{Riemann tensor}{85}
\idx{0}{$SL(2,\mathbb C)$（代数）}{$SL(2,\mathbb C)$ (algebra)}{71}
\idx{1}{不変状態}{invariant state}{71, 203}
\idx{0}{$SL(2,\mathbb R)$（代数）}{$SL(2,\mathbb R)$ (algebra)}{56, 71}
\idx{0}{$SO(D-2)$}{}{23–4, 27, \textit{30}}
\idx{0}{$SO(D-1)$}{}{23–4, 27, \textit{30–1}}
\idx{0}{$SO(n)$}{}{30, 191–2, 228}
\idx{1}{D-brane 上のもの}{on D-brane}{279–80}
\idx{0}{$SO(8192)$}{}{228–9, 280}
\idx{0}{$Sp(k)$}{}{30, 191–2, 228, 279}
\idx{0}{$SU(n)$}{}{189}
% Original index: p.400
\idx{0}{$SU(2)\times SU(2)$ 半径}{$SU(2)\times SU(2)$ radius}{242–8, 261–3}
\idx{1}{D-brane 上のもの}{on D-brane}{264–5, 269–72, 279}
\idx{0}{$SU(3)\times SU(2)\times U(1)$ 模型}{$SU(3)\times SU(2)\times U(1)$ model}{1}
\idx{1}{}{}{標準模型（Standard Model）も参照}
\idx{0}{$S$ 行列}{$S$-matrix}{97–101, 103–4}
\idx{1}{計算}{calculations}{178–205}
\idx{1}{その表示}{expression for}{99, 155–7, 164}
\idx{1}{ユニタリティ}{unitarity}{181–2, 186–7, 193–4, 196, 285–90, 306}
\idx{1}{Green 関数との比較}{vs Green's function}{104}
\idx{0}{$s,t,u$（Mandelstam 変数）}{$s,t,u$ (Mandelstam variables)}{180}
\idx{0}{飽和する測地線}{saturating geodesic}{300}
\idx{0}{スケール変換}{scale transformation}{44–8, 90–1}
\idx{0}{Schottky 群}{Schottky group}{290–1, 351}
\idx{0}{Schwarz 微分}{Schwarzian}{48}
\idx{0}{Schwinger–Dyson 方程式}{Schwinger–Dyson equation}{334}
\idx{0}{自己縮約}{self-contractions}{103, 171, 338}
\idx{0}{自己双対半径}{self-dual radius}{242–8, 261–3}
\idx{0}{縫い合わせ}{sewing}{332, 351}
\idx{1}{円}{circles}{298–9}
\idx{1}{3 点を穿孔した球面からの一般の曲面の構成}{construction of general surface from three-punctured spheres}{296–7}
\idx{1}{CFT の縫い合わせ}{of CFTs}{300–5}
\idx{2}{共形ブートストラップとの関係}{and conformal bootstrap}{303–5}
\idx{1}{Riemann 面の縫い合わせ}{of Riemann surfaces}{294–300}
\idx{1}{パラメーター}{parameter}{294–5}
\idx{1}{弦振幅の解析への利用}{used to analyze string amplitudes}{305–7}
\idx{0}{Siegel ゲージ}{Siegel gauge}{313}
\idx{0}{シグマ模型計量}{sigma model metric}{115}
\idx{0}{ソリトン}{soliton}{104, 235}
\idx{0}{球面}{sphere}{}
\idx{1}{期待値}{expectation values}{174–8, 198–204}
\idx{1}{Riemann 面}{Riemann surface}{166–8}
\idx{1}{弦振幅}{string amplitudes}{192–6, \textit{204–5}, 283–90}
\idx{1}{3 点を穿孔したもの}{thrice-punctured}{296–7, 303}
\idx{2}{射影平面との相殺}{cancellation with projective plane}{229}
\idx{0}{スピン}{spin}{23–4, 184, \textit{204}}
\idx{1}{CFT におけるもの}{in CFT}{46, 213}
\idx{1}{Coleman–Weinberg 公式におけるもの}{in Coleman–Weinberg formula}{221}
\idx{0}{偽状態}{spurious state}{122}
\idx{0}{標準模型}{Standard Model}{1–7, 25, 112, 310, 348–9}
\idx{0}{状態・演算子同型}{state–operator isomorphism}{63–73, 101–2}
\idx{1}{具体形}{explicit form}{65}
\idx{1}{経路積分による導出}{path integral derivation}{66–8}
\idx{1}{頂点演算子}{vertex operators}{101–2, 164, 199–201}
\idx{0}{Stirling 近似}{Stirling's approximation}{184, 315}
\idx{0}{弦の結合定数}{string coupling}{83–4}
\idx{1}{ディラトンとの関係}{and dilaton}{112–13}
\idx{1}{作用の全体因子として}{as overall factor in action}{115}
\idx{1}{頂点演算子中のもの}{in vertex operators}{101–2, 107, 179}
\idx{1}{位置に依存するもの}{position-dependent}{323–4}
\idx{1}{開弦と閉弦の間の関係}{relation between open and closed}{83–4, 276}
\idx{1}{ゲージ結合との関係}{relation to gauge coupling}{187, 234}
\idx{1}{重力結合との関係}{relation to gravitational coupling}{194, \textit{205}}
\idx{0}{弦の場の理論}{string field theory}{310–17, 352}
\idx{0}{弦の長さスケール}{string length scale}{25, 110, 244, 248}
\idx{0}{弦計量}{string metric}{115}
\idx{0}{弦の張力}{string tension}{11, \textit{30}}
\idx{0}{減算}{subtraction}{39, 47, 60}
\idx{0}{超対称性}{supersymmetry}{2, 6–7, 222, 283, 315, 348}
\idx{0}{超対称性の破れ}{supersymmetry breaking}{222, 233–4}
\idx{0}{T 双対性}{T-duality}{247–8, \textit{280–1}}
\idx{1}{D-brane との関係}{and D-branes}{268–73}
\idx{1}{ゲージ対称性として}{as gauge symmetry}{248}
\idx{1}{$O(k,k,\mathbb Z)$}{}{253–5}
\idx{1}{開弦理論}{open string theory}{265–8}
\idx{1}{巻き付きと運動量の交換}{reverses winding and momentum}{247, 268}
\idx{1}{向きのない弦理論}{unoriented string theory}{277–80}
\idx{0}{タキオン}{tachyon}{}
\idx{1}{振幅}{amplitudes}{178–84, 187, 192–5, \textit{204–5}, 217, \textit{229–30}}
\idx{1}{ボソン弦理論に起因するもの}{artifact of bosonic string theory}{23, 219, 223–4, 246, 275, 306}
\idx{1}{BRST 量子化におけるもの}{in BRST quantization}{135}
\idx{1}{光円錐量子化におけるもの}{in light-cone quantization}{23, 27, 29–30}
\idx{1}{線形ディラトン背景におけるもの}{in linear dilaton background}{323–4, \textit{328}}
\idx{1}{OCQ におけるもの}{in OCQ}{123}
% Original index: p.401
\idx{1}{すべてのボソン弦理論に存在}{present in all bosonic string theories}{117}
\idx{1}{頂点演算子}{vertex operator}{101, 107}
\idx{0}{タドポール}{tadpole}{197, 309}
\idx{1}{$SO(8192)$ における相殺}{cancellation for $SO(8192)$}{229}
\idx{1}{円板}{disk}{197, 223–6, 228–9}
\idx{1}{射影平面}{projective plane}{197, 227–9}
\idx{1}{トーラス}{torus}{307}
\idx{0}{標的空間}{target space}{109}
\idx{0}{Teichmüller パラメーター}{Teichmüller parameters}{148}
\idx{0}{温度}{temperature}{320–2, \textit{328}}
\idx{0}{テンソル演算子（プライマリー場）}{tensor operator (primary field)}{46–7, \textit{76}}
\idx{1}{期待値}{expectation values}{199}
\idx{1}{最高ウェイト状態への写像}{maps to highest weight state}{71}
\idx{1}{OPE}{}{47, 200, 303}
\idx{0}{標構}{tetrad}{}
\idx{1}{時空}{spacetime}{250–1}
\idx{1}{世界線}{world-line}{10, 129, 145–6}
\idx{1}{世界面}{world-sheet}{83}
\idx{0}{テータ関数}{theta functions}{214–16, \textit{229}, 350}
\idx{1}{期待値中のもの}{in expectation values}{208–9, 211, 228, 259}
\idx{0}{トーラスコンパクト化}{toroidal compactification}{231–55, 236–80, \textit{280–1}, 351}
\idx{1}{抽象的（Narain）記述}{abstract (Narain) description}{251–4}
\idx{1}{CFT におけるもの}{in CFT}{235–40}
\idx{1}{場の理論におけるもの}{in field theory}{231–5}
\idx{1}{弦理論におけるもの}{in string theory}{241–8}
\idx{1}{複数の次元のもの}{of several dimensions}{249–55}
\idx{1}{2 次元のもの}{of two dimensions}{255}
\idx{1}{開弦}{open string}{263–8}
\idx{0}{トーラス}{torus}{}
\idx{1}{分配関数との関係}{and partition function}{209}
\idx{1}{期待値}{expectation values}{208–13}
\idx{2}{1 点関数}{one-point function}{304}
\idx{1}{ねじれた CFT におけるもの}{in twisted CFT}{260}
\idx{1}{Riemann 面}{Riemann surface}{146–50, 154, 206–7}
\idx{1}{弦振幅}{string amplitude}{216–22, \textit{229–30}, 237, 250}
\idx{2}{有限性}{finiteness}{218–20}
\idx{0}{トレース}{trace}{}
\idx{1}{エネルギー運動量テンソルのもの}{of energy-momentum tensor}{13, 43, 92–7, 111}
\idx{1}{分配関数}{partition function}{209, 335–6, 344, \textit{345–6}}
\idx{0}{遷移関数}{transition functions}{153–4, 161–4}
\idx{1}{向きのない場合}{unoriented}{277}
\idx{0}{横方向}{transverse}{}
\idx{1}{Hilbert 空間}{Hilbert space}{138}
\idx{1}{振動子}{oscillators}{20–1, 26, 125, 217–18}
\idx{0}{木レベル振幅}{tree-level amplitudes}{178–98, 283–90, 350–1}
\idx{1}{継承原理}{inheritance principle}{258, 261}
\idx{1}{開弦と閉弦の関係}{open–closed relation}{195, 350}
\idx{1}{極}{poles}{181–2, 193–4, 196, 287–8}
\idx{0}{ねじれた場}{twisted fields}{\textit{31}, 257, 259, 336, \textit{345–6}}
\idx{0}{ねじれた状態}{twisted state}{\textit{31}, 257–61}
\idx{0}{ねじり（オービフォールド）構成}{twisting (orbifold) construction}{259–63}
\idx{0}{$U(n)$}{}{30, 185, 188–9, \textit{230}}
\idx{0}{紫外（UV）発散}{ultraviolet (UV) divergences}{発散（divergences）の項を参照}
\idx{0}{単位ゲージ}{unit gauge}{85}
\idx{0}{ユニタリティ}{unitarity}{$S$ 行列（$S$-matrix）の項を参照}
\idx{0}{ユニタリー CFT}{unitary CFT}{71–3, 117, 137, 213}
\idx{0}{向きのない弦理論}{unoriented string theory}{29–30, 80–1}
\idx{1}{Chan–Paton 状態}{Chan–Paton states}{189–92}
\idx{1}{1 ループ振幅}{one-loop amplitudes}{226–9}
\idx{1}{ゲージ対称性としての向きの反転}{orientation-reversal as gauge symmetry}{149, 153, 277–9}
\idx{1}{Riemann 面}{Riemann surfaces}{80–1, 100–1}
\idx{1}{T 双対性}{T-duality}{277–80, 351}
\idx{0}{真空振幅}{vacuum amplitude}{}
\idx{1}{場の理論}{field theory}{218}
\idx{1}{その物理}{physics of}{220–2}
\idx{1}{弦}{string}{198, 216–30, 274–6}
\idx{0}{真空の縮退}{vacuum degeneracy}{7, 118, 233–4, 261, 348}
\idx{0}{Veneziano 振幅}{Veneziano amplitude}{178–84, \textit{204}, 348}
\idx{0}{頂点演算子}{vertex operators}{63–8, \textit{74}, 97–102, 349}
\idx{1}{一般的な構成}{general construction}{63–6}
\idx{1}{質量殻上になければならないこと}{must be on mass-shell}{101–3, 113}
\idx{1}{規格化}{normalization}{101–2, 107, 187, 194, 276, 284}
\idx{1}{OCQ}{}{125, 143, 160–1, 164, \textit{165}}
\idx{1}{Polyakov の方法}{Polyakov approach}{102–7, \textit{119}}
\idx{1}{位置}{positions}{99–100}
\idx{2}{一致極限}{coincident limit}{195–6, 306–9}
\idx{2}{具体的な振幅におけるもの}{in specific amplitudes}{178–80, 192–3, \textit{204}, 216–17}
% Original index: p.402
\idx{2}{モジュライおよび CKV との関係}{relation to moduli and CKVs}{149–52, 154–7, 163–4}
\idx{1}{具体例}{specific}{101–2, 107}
\idx{1}{ねじれた状態}{twisted state}{258, 351}
\idx{1}{巻き付き状態}{winding state}{238–9, 243–5, 251}
\idx{0}{Virasoro 代数}{Virasoro algebra}{52–7, 75–6}
\idx{0}{Virasoro 生成子}{Virasoro generators}{}
\idx{1}{拘束条件として}{as constraints}{121–2, 133}
\idx{1}{演算子を中心とするもの}{centered on operators}{70}
\idx{2}{微分としての $L_{-1}\cdot$}{$L_{-1}\cdot$ as derivative}{71}
\idx{2}{ウェイトとしての $L_0\cdot$}{$L_0\cdot$ as weight}{71}
\idx{1}{Hermite 性}{Hermiticity}{71}
\idx{1}{振動子表示}{oscillator expressions}{59, 61–2}
\idx{1}{巻き付き状態}{winding state}{238}
\idx{0}{Virasoro–Shapiro 振幅}{Virasoro–Shapiro amplitude}{192–5, \textit{204–5}}
\idx{0}{Ward 恒等式}{Ward identity}{41–3, 55}
\idx{1}{局所的な性質として}{as local property}{86}
\idx{1}{共形対称性に対するもの}{for conformal symmetry}{45–6}
\idx{0}{ウェイト（共形）}{weight (conformal)}{46–7, 71–3}
\idx{1}{頂点演算子のもの}{of vertex operators}{101, 125, 143}
\idx{1}{正値性}{positivity}{72}
\idx{0}{Wess–Zumino 整合性条件}{Wess–Zumino consistency condition}{96}
\idx{0}{Weyl 異常}{Weyl anomaly}{異常（anomaly）の項を参照}
\idx{0}{Weyl 変換}{Weyl transformation}{13, 15–16, \textit{30}}
\idx{1}{冗長性として}{as redundancy}{13, 84, 117, 326}
\idx{1}{そのゲージ固定}{gauge-fixing of}{17–18, 85}
\idx{1}{ゴーストは自明}{ghosts trivial}{88, 131}
\idx{1}{曲率の変換}{of curvature}{15, 85, 114}
\idx{1}{非線形シグマ模型の変換}{of nonlinear sigma model}{110–12}
\idx{1}{頂点演算子の変換}{of vertex operators}{102–7}
\idx{1}{共形変換との関係}{relation to conformal transformation}{86}
\idx{1}{時空}{spacetime}{114, 233, 254}
\idx{0}{Wick の定理}{Wick's theorem}{60}
\idx{0}{Wilson 線}{Wilson line}{263–8, 279–80}
\idx{1}{ゲージ対称性を破ること}{breaks gauge symmetry}{265, 270}
\idx{0}{巻き付き状態}{winding state}{\textit{31}, 235–9}
\idx{1}{Dirichlet 開弦}{Dirichlet open string}{266–7}
\idx{1}{$R\to0$ で質量ゼロ}{massless as $R\to0$}{247}
\idx{1}{頂点演算子}{vertex operators}{238–9, 243–5, 251}
\idx{0}{世界面}{world-sheet}{13–14}
\idx{0}{世界面双対性}{world-sheet duality}{287, 332}
\idx{1}{円筒}{cylinder}{226, 274}
\idx{0}{巻き付いたもの}{wrapped}{}
\idx{1}{D-brane}{}{273}
\idx{1}{世界面}{world-sheet}{237, 250}
\idx{0}{Yukawa 結合}{Yukawa couplings}{1, 79}
\idx{0}{Zamolodchikov 内積}{Zamolodchikov inner product}{200}
\idx{0}{ゼロモード}{zero modes}{}
\idx{1}{ゴースト}{ghost}{157–9, 161, 176, 212}
\idx{1}{$X^\mu$}{}{169–70, 173}
\idx{0}{零ノルム状態}{zero-norm states}{零状態（null states）を参照}
\idx{0}{零点エネルギー}{zero-point energy}{59–61, 72}
\idx{1}{発見的な議論}{heuristic}{21–2, 26}
\idx{1}{覚え方}{mnemonics}{73, 125}
\idx{0}{零勾配極限}{zero-slope limit}{188}
\endgroup
